% Related work (Task 42): the submitted text, trimmed; every citation of the submitted version kept; positioning
% paragraph updated for A8x (handoff/37_narrative_v2.md section 6, rows 2 and 3).
\section{Related Work}
\label{sec:related}

\paragraph{Attention as explanation.}
Whether attention weights explain model decisions has been argued in both directions for text models
\cite{jain2019,wiegreffe2019,serrano2019}, and Jacovi and Goldberg argue for a graded notion of faithfulness and give
guidelines for its evaluation \cite{jacovi2020}.
Because a Transformer mixes information across layers, attention rollout and attention flow aggregate attention across
layers \cite{abnar2020}, relevance propagation combines attention with gradients \cite{chefer2021}, and ALTI also
accounts for the residual connection and layer normalisation \cite{ferrando2022}.
For tabular data, TabNet derives feature attributions from its sequential attention \cite{arik2021}, the authors of the
FT-Transformer averaged the attention maps of the classification token ([CLS]) as a feature-importance estimate
\cite{gorishniy2021}, and attention maps of tabular transformers have been used qualitatively \cite{somepalli2021};
neither use establishes that the maps are faithful.

\paragraph{Attribution methods and faithfulness evaluation.}
We compared attention with Integrated Gradients \cite{sundararajan2017}, GradientSHAP, which builds on Shapley additive explanations (SHAP) \cite{lundberg2017} and is computed as expected gradients \cite{erion2021},
occlusion \cite{zeiler2014}, and TreeSHAP on a gradient-boosted reference model \cite{lundberg2020}, using deletion and
insertion curves \cite{petsiuk2018,samek2017}, the perturbation curve gap \cite{samek2017} and a parameter-randomisation
check \cite{adebayo2018}.
Removal-based scores can be confounded by the distribution shift of the masked inputs \cite{hooker2019}, can rank
methods inconsistently across evaluation strategies \cite{rong2022}, and can be statistically unreliable
\cite{tomsett2020}; our curves mask with the training mean without retraining and share these caveats.

\paragraph{Fragility of explanations and of explanation-based monitors.}
Explanations can be changed by small input perturbations without changing the prediction
\cite{ghorbani2019,dombrowski2019}, attention can be made deceptive through training \cite{pruthi2020}, and
explanations can be manipulated through the model \cite{heo2019,slack2020}; Baniecki and Biecek survey these attacks
and the proposed countermeasures \cite{baniecki2024}, and a review of user-centric intelligent systems treats
explainability together with fairness \cite{zahid2025}.
Explanation methods have been used to detect adversarial inputs \cite{fidel2020}, and Zhang et al. construct inputs that
mislead a classifier and its interpretation model at the same time \cite{zhang2020}.
Detectors of adversarial inputs have repeatedly been bypassed by attacks designed for them, so such detectors need to be
evaluated with adaptive attacks that model the detector \cite{carlini2017,tramer2020}.

\paragraph{Explainability for network intrusion detection.}
Warnecke et al. propose criteria for evaluating explanation methods in security \cite{warnecke2020}; Tritscher et al.
find that the quality of SHAP explanations for NIDS depends on the replacement data \cite{tritscher2023}; and Arreche
and Abdallah compare white-box and black-box methods on three intrusion datasets, among them CIC-IDS2017 (from the
Canadian Institute for Cybersecurity, CIC) \cite{arreche2025}.
Several NIDS studies pair ensemble detectors with explanations to describe the detector or to select its features
\cite{ahmed2024haenid,ahmed2025hcivad,ahmed2025consensus,ahmed2024shap}, deep learning has been applied to
denial-of-service (DoS) and distributed denial-of-service (DDoS) detection in 5G networks \cite{albashayreh2025}, and
transformer models for flow-based NIDS \cite{manocchio2024} and deep models for tabular data \cite{borisov2024} have been
compared.
Closest to our monitoring question, Maseno et al. attack recurrent intrusion detection system (IDS) models with the fast
gradient sign method (FGSM) and projected gradient descent (PGD) and use the drift of SHAP explanations as an auxiliary
detection signal, which they report as informative but not perfectly separable at larger scale \cite{maseno2026}.

\paragraph{Adversarial attacks on NIDS and dataset quality.}
Feature-space attacks on NIDS can produce flows that no real traffic could generate, which motivates problem-space
constraints and validity checks \cite{pierazzi2020,alhussien2024} and realistic threat models \cite{apruzzese2022};
constrained attacks on NIDS include JSMA under domain constraints \cite{sheatsley2022} and black-box botnet evasion
\cite{debicha2023}, and Khaleel et al. review countermeasures \cite{khaleel2024}.
Our constraint model follows this line in feature space; it does not generate packets.
LYCOS-IDS2017 re-extracts the features of CIC-IDS2017 after an analysis of problems in the original flows
\cite{rosay2021,rosay2022}, and LycoS-Unicas-IDS2018 applies the corresponding corrections to CSE-CIC-IDS2018 (a collaboration of the Communications Security Establishment, CSE, and CIC) \cite{cantone2024}.
Issues of the original datasets \cite{sharafaldin2018} are documented in \cite{engelen2021,liu2022,lanvin2022}, of
benchmark NIDS datasets in general in \cite{flood2024}, and pitfalls of learning-based security evaluations in
\cite{arp2022}.

\paragraph{Positioning.}
These lines of work treat the faithfulness of attention, explanation-based monitors and constrained attacks on NIDS
largely separately.
We bring them together for one model and task: an FT-Transformer on two corrected flow datasets, with attention
compared to four attribution methods on clean data and under protocol-constrained attacks, and attention-attribution
consistency scores evaluated as an addition to a confidence-and-distance detector against an attacker that ignores them,
one that models them through differentiable stand-ins and one that optimises their exact values, with control detectors
that separate the role of attention.
We do not propose a new explanation method or a new detector, and we do not claim that attention is unfaithful in
general; the consistency scores gave no gain against the strongest adaptive attacker we built.
